\section{Discussion}
\label{sec:discussion}

The theorem suite gives two complementary normal forms for the convex-to-polytope conjecture.

The first is \emph{local and rational}.  Universal polyhedralization reduces to binary-meet deletion.  Every such deletion has a finite rational interval bridge contained in the strict trunk above the missing word.  If an appropriate selector exists over a fixed polytopal lower realization, finite cell separation makes the selector polytopal in the same dimension.  Hence neither a continuum of bridge events nor curvature of one selector over a fixed arrangement is the obstruction.

The second is \emph{global and finite}.  Every source realization has a protected inner-polytopal supercode with exactly the same maximal words and a finite set of unwanted atoms.  Every complete unwanted carrier has a finite cover by source-contained polytopes labelled by complete target words.  Thus degeneracy and zero margin do not force an infinite collection of local charts.  The obstruction is the simultaneous convex globalization of repeated neuron labels across those finitely many modules.

\subsection{Oriented-matroid interpretation}

If the remaining globalization theorem is proved, the final receptive fields have finitely many facets.  Take one affine hyperplane for every facet of every neuron polytope and of the ambient polytope.  The resulting finite real hyperplane arrangement represents an oriented matroid $M$.  Each original neuron is the intersection of the positive halfspaces corresponding to its facets, hence is represented by a trunk of $L^+(M)$.  Restricting to the ambient trunk and reading those neuron trunks gives a surjective code morphism
\[
 L^+(M)\twoheadrightarrow\C.
\]
This is the constructive content behind the characterization of \citet{KuninLienkaemperRosen2023}.

Conversely, a negative answer would require more than the failure of a chosen approximation, cap, bridge, or finite lift.  One would need an explicit convex code whose every finite representable completion fails.  None of the exact failure examples in this paper has that invariance: each obstructs a particular globalization rule, not polytope convexity of the code.

\subsection{The remaining theorem}

The unresolved statement can be formulated without any limiting language.

\begin{quote}
\textbf{Finite simultaneous attachment problem.}
Given a protected inner-polytopal supercode $\D$ of an open convex code $\C$ and one finite source-contained whole-word repair atlas for every word in $\D\setminus\C$, construct finitely many polytopes, possibly in a larger dimension, that activate each atlas label synchronously, eliminate every unwanted carrier, preserve every protected target word, and represent every repeated original neuron by one convex polytope.
\end{quote}

A solution would imply universal polyhedralization directly.  By \Cref{thm:binary-meet-reduction}, it would also solve the canonical binary-meet deletion operation.  The rational bridge theorem supplies the one-dimensional interface data required in that local formulation.

\subsection{Algorithmic scope}

All objects constructed in the present reductions are finite once an input realization and atom witnesses are given.  This does not by itself yield a decision algorithm for convexity: the input realization is not combinatorial data, and no computable universal bound on a final completion is proved.  Likewise, a bounded search for oriented-matroid completions cannot certify nonpolytopality without a separate completion-size bound.

\subsection{Biological interpretation}

The code records exact co-firing regions rather than only intersections of receptive fields.  Replacing smooth place fields by polytopes would show that the finite combinatorics of co-firing never requires curved receptive-field boundaries.  The present results establish this in the plane and for the finite classes above, and show that the unresolved issue is how several appearances of the same receptive field couple across a finite collection of local repair modules.  We make no claim that biological receptive fields are polytopal, only that their finite atom codes may be.

In summary, the conjecture has been reduced to a canonical binary-meet deletion with finite rational bridge data and, independently, to a finite simultaneous globalization problem for protected whole-word repair atlases.  These reductions eliminate uncontrolled local curvature and infinite local repair as possible explanations for the difficulty, but they do not settle the conjecture.
